% !PRINTED 213
\weberpage{213}
To find the remaining substitutions of the group, we observe that by the theorem \S.\,53, 2., if we take $\pm 1$ for $a, a'$ and $0, \infty$ for $b, b'$, we can infer the existence of a substitution $\chi$ in $G$ which satisfies the conditions
\[
\chi(-1) = 0, \quad \chi(+1) = \infty,
\]
and that $\chi$ must therefore have the form:
\weq{2}{\chi = \lambda\,\frac{x+1}{x-1},}
where $\lambda$ is a constant factor.

If we substitute the six conjugate poles $0, \infty, +1, -1,$
$+i, -i$ into $\chi$, we must obtain the same values in another order (\S.\,52, 5.). But these values are
\[
-\lambda,\ \lambda,\ \infty,\ 0,\ -\lambda i,\ \lambda i.
\]

Hence $\lambda = \pm 1$ or $= \pm i$ must hold. The values $\pm 1$ are not admissible, since otherwise
\[
\chi(i) = \mp i, \quad \Theta\chi(i) = \pm i,
\]
and so $\pm i$ would be the poles of $\chi$ or of $\Theta\chi$, whereas they are in fact only two-fold and are the poles of $\Theta\psi$. Hence $\lambda = \pm i$ must hold, and we may take the upper sign without restriction of generality, for the lower sign then corresponds to the substitution $\Theta\chi$. Therefore
\[
\chi(x) = i\,\frac{x+1}{x-1}.
\]

From this one obtains
\[
\chi^2(x) = \frac{x+i}{x-i}, \quad \chi^3(x) = x,
\]
and there result the 12 substitutions:
\weq{3}{\begin{array}{cccc}
1, & \Theta, & \psi, & \Theta\psi \\
\chi, & \Theta\chi, & \psi\chi, & \Theta\psi\chi \\
\chi^2, & \Theta\chi^2, & \psi\chi^2, & \Theta\psi\chi^2,
\end{array}}
or
\weq{4}{\pm x,\ \pm\frac{1}{x},\ \pm i\,\frac{x+1}{x-1},\ \pm i\,\frac{x-1}{x+1},\ \pm\frac{x+i}{x-i},\ \pm\frac{x-i}{x+i}.}

That these substitutions really form a group, the \gesperrt{tetrahedral group}, of which $\Theta, \psi, \chi$ are the generating elements, follows from
\weq{5}{\chi\Theta = \Theta\psi\chi,\ \chi^2\Theta = \psi\chi^2,\ \psi\Theta = \Theta\psi,\ \chi\psi = \Theta\chi,\ \chi^2\psi = \Theta\psi\chi^2,}
%CONT>>
